\documentclass[12pt]{article}
\usepackage[utf8]{inputenc}
\usepackage[T1]{fontenc}
\usepackage{hyperref}
\title{Fali of Baissa Language Resources for K-12 Curriculum Development}
\author{Compiled by Claude Sonnet 5.5 (claude-sonnet-5-5) for Andr\'as Kornai}
\date{}
\begin{document}
\maketitle

\section{Introduction}
Fali of Baissa (ISO 639-3: fah; Glottocode bais1242) is an unwritten Benue-Congo language once spoken at Baissa, Taraba State, Nigeria. Wikipedia says it went extinct in the 2000s. Roger Blench's \emph{Atlas of Nigerian Languages} (2012), as mirrored by the Endangered Languages Project, says it was still spoken by a few people on the Falinga Plateau; the ELP rates it critically endangered with 1--9 speakers, at only 20 percent certainty. Its status after 2012 is unconfirmed.

\section{Primary Educational Resources}
None found.

\section{Linguistic and Reference Materials}
\begin{itemize}
\item Wikipedia, Fali of Baissa. \url{https://en.wikipedia.org/wiki/Fali_of_Baissa}
\item Endangered Languages Project page, with information from Blench (2012). \url{https://endangeredlanguages.com/lang/4144?hl=en}
\item Endangered Languages Project bibliography for the language. \url{https://endangeredlanguages.com/lang/4144/bibliography}
\end{itemize}

\section{Religious and Community Resources}
None found.

\section{Digital and Audio Resources}
None found in search results.

\section{Curriculum Development Context}
Almost nothing exists. Any heritage curriculum would have to rest on new documentation, and on confirming whether any speakers remain. It could draw on the bibliography of earlier work listed by the Endangered Languages Project.

\end{document}